\documentclass[aps,prc,onecolumn,superscriptaddress,nofootinbib,10pt]{revtex4-2}
\usepackage{amsmath,amssymb}
\usepackage{booktabs}
\usepackage[colorlinks=true,linkcolor=blue,citecolor=blue,urlcolor=blue]{hyperref}

\newcommand{\nB}{n_B}
\newcommand{\nC}{n_C}
\newcommand{\nS}{n_S}
\newcommand{\muB}{\mu_B}
\newcommand{\muC}{\mu_C}
\newcommand{\muS}{\mu_S}
\newcommand{\hc}{(\hbar c)^3}
\newcommand{\nbs}{\bar{n}^s}
\newcommand{\SR}{\Sigma^{R}}
\newcommand{\kF}{k_F}

\begin{document}

\title{ENJL: the extended Nambu--Jona-Lasinio model of baryonic and quark
       matter \\
       --- as implemented in \texttt{eos/enjl}}

\author{M.~Guerrini}
\affiliation{University of Ferrara}

\begin{abstract}
The \texttt{eos/enjl} subpackage implements the extended Nambu--Jona-Lasinio
model of Xia~\cite{Xia2024}, in which baryons and quarks are described by
\emph{one} functional: a baryon is a three-quark cluster whose mass is built
from the same constituent quark masses that the NJL gap equation
determines, so the chiral, quarkyonic and deconfinement transitions all come
out of a single mean field rather than out of two models joined at a
boundary. The independent unknowns of that mean field are the three
constituent masses $M_u$, $M_d$, $M_s$; everything else --- the baryon masses,
the vector fields, the rearrangement terms --- follows algebraically. The
model carries density-dependent couplings $\alpha_S(\nB)$, $\Gamma_\omega(\nB)$
and $\Gamma_\rho(\nB)$, and therefore rearrangement self-energies, and it
carries the NJL three-momentum cut-off $\Lambda$, which enters as a
temperature-independent vacuum subtraction on the quark sector alone.

This note states the model as the code solves it: the single-species
thermodynamics at $T = 0$ and at $T > 0$ with the cut-off written out, the
gap equation in
the form the code uses rather than the form the paper prints, the vector and
rearrangement sector, the assembly of $\varepsilon$, $P$, $s$ and $n^s$, the
conserved-charge bookkeeping in this repository's conventions, the
equilibrium residual row by row, and how each of the repository's four modes
--- all four of which are closed --- is assembled from one declaration rather
than from four solvers. Section~\ref{sec:validation}
quotes what the implementation actually reproduces, measured against the
paper and against the author's own tables. Reference~\cite{Xia2024} is a
$T = 0$ paper; the implementation is closed at any non-negative temperature,
with entropy per baryon accepted wherever a temperature is, and
Sec.~\ref{sec:numerics} says what temperature costs and what is still open
above the first first-order transition.
\end{abstract}

\maketitle

\section{Degrees of freedom}
\label{sec:dof}

Eight species, in three families. The quantum numbers are those of
\texttt{eos.enjl.species}, in this repository's conventions (CLAUDE.md
\S2): $C$ is the electric charge of \emph{strongly interacting} matter only,
and strangeness is $S = +1$ per $s$ quark --- so $\Lambda$ has $S = +1$, the
opposite of the PDG sign.

\begin{center}
\begin{tabular}{lcccccl}
\toprule
species & $g$ & $B$ & $C$ & $S$ & $\tau_3$ & valence content
                                              $(N^u_i, N^d_i, N^s_i)$ \\
\midrule
$p$       & 2 & $1$   & $+1$   & $0$  & $+1$ & $(2,1,0)$ \\
$n$       & 2 & $1$   & $0$    & $0$  & $-1$ & $(1,2,0)$ \\
$\Lambda$ & 2 & $1$   & $0$    & $+1$ & $0$  & $(1,1,1)$ \\
\midrule
$u$       & 6 & $1/3$ & $+2/3$ & $0$  & $+1$ & $(1,0,0)$ \\
$d$       & 6 & $1/3$ & $-1/3$ & $0$  & $-1$ & $(0,1,0)$ \\
$s$       & 6 & $1/3$ & $-1/3$ & $+1$ & $0$  & $(0,0,1)$ \\
\midrule
$e^-$     & 2 & $0$   & ---    & $0$  & $0$  & --- \\
$\mu^-$   & 2 & $0$   & ---    & $0$  & $0$  & --- \\
\bottomrule
\end{tabular}
\end{center}

\noindent
The quark degeneracy $g = 6$ is three colours times two spins; the baryon and
lepton degeneracy $g = 2$ is spin alone. The leptons carry electric charge
$q_e = q_\mu = -1$ but are excluded from $C$ by the convention above; total
electric neutrality is imposed separately where a mode calls for it
(Sec.~\ref{sec:beta}). The last column, $N^q_i$, is the number of valence
quarks of flavour $q$ in baryon $i$, and it is the object through which the
whole model is built: it appears in the baryon mass
Eq.~\eqref{eq:baryonmass}, in the effective scalar density
Eq.~\eqref{eq:nbar}, in the $\omega$ source Eq.~\eqref{eq:Jomega}, and in the
rearrangement term Eq.~\eqref{eq:SigmaRb}.

A model-wide rescaling factor $f_i$ multiplies each species' coupling to the
vector fields: $f_p = f_n = 1$ by construction, $f_\Lambda$ is fitted to the
$\Lambda$ potential depth, and the three quarks share one value $f_q$ which is
a parameter of the study (Sec.~\ref{sec:parameters}). Leptons do not couple,
$f_e = f_\mu = 0$.

\section{Single-species thermodynamics, and the cut-off}
\label{sec:kinetic}

Each species is a free Fermi gas of degeneracy $g$ and effective mass $M$ at
kinetic (effective) chemical potential $\nu$, and it is the \emph{only} place
temperature enters. At $T = 0$ the gas is filled to a sharp Fermi momentum
$\kF = \sqrt{\nu^2 - M^2}$ and every integral below has a closed form; at
$T > 0$ there is no sharp surface, $\kF$ is not defined, and $\nu$ is the
primary variable. Section~\ref{sec:kineticT} gives the finite-temperature
integrals and Sec.~\ref{sec:assembly} the entropy they carry.

Writing $x = \kF/M$ and $\nu = \sqrt{\kF^2 + M^2}$ for the
kinetic (effective) chemical potential, and $\hc = 7.6838\times10^{6}\;
\mathrm{MeV^3\,fm^{-3}}$ for the conversion to the $\mathrm{fm^{-3}}$ and
$\mathrm{MeV\,fm^{-3}}$ used at every public boundary,
%
\begin{align}
n^{\rm med}(\kF, M, g) &= \frac{g\,\kF^3}{6\pi^2} ,
\label{eq:nmed}\\[2pt]
n^{s,{\rm med}}(\kF, M, g) &= \frac{g M^3}{4\pi^2}
   \Big[\, x\sqrt{x^2+1} - \operatorname{arcsinh} x \,\Big] ,
\label{eq:nsmed}\\[2pt]
\varepsilon^{\rm med}(\kF, M, g) &= \frac{g M^4}{16\pi^2}
   \Big[\, x(2x^2+1)\sqrt{x^2+1} - \operatorname{arcsinh} x \,\Big] ,
\label{eq:epsmed}\\[2pt]
P^{\rm med}(\kF, M, g) &= \nu\, n^{\rm med} - \varepsilon^{\rm med}
 = \frac{g M^4}{48\pi^2}
   \Big[\, x(2x^2-3)\sqrt{x^2+1} + 3\operatorname{arcsinh} x \,\Big] ,
\label{eq:Pmed}\\[2pt]
s &= 0 \qquad\text{(every species, at $T = 0$)} .
\label{eq:s}
\end{align}
%
Equation~\eqref{eq:epsmed} includes the rest mass, and $\varepsilon + P =
\nu n$ holds identically for one species by Eq.~\eqref{eq:Pmed}. Below
threshold, $\nu \le M$, all four vanish. These are the standard
$T = 0$ Fermi integrals and are \emph{not} written a second time inside this
model: the code calls
\texttt{eos.general.fermi\_integrals.solve\_fermi\_jel} at $T = 0$, which
routes to the exact closed forms above. They are restated here because a
paper-style description must be self-contained.

\subsection{Finite temperature}
\label{sec:kineticT}

At $T > 0$ the occupation is Fermi--Dirac rather than a step, and
antiparticles are thermally populated. With $E(k) = \sqrt{k^2 + M^2}$ and
%
\begin{equation}
f_{\mp}(k) \;=\; \frac{1}{1 + \exp\!\big[\,(E(k) \mp \nu)/T\,\big]}
\label{eq:fermidirac}
\end{equation}
%
the occupation of particles ($f_-$) and antiparticles ($f_+$), the four
medium integrals of Eqs.~\eqref{eq:nmed}--\eqref{eq:Pmed} become
%
\begin{align}
n^{\rm med}(\nu, M, g, T) &= \frac{g}{2\pi^2}\int_0^\infty\!\! dk\;k^2
   \big[\, f_-(k) - f_+(k) \,\big] ,
\label{eq:nmedT}\\[2pt]
n^{s,{\rm med}}(\nu, M, g, T) &= \frac{g}{2\pi^2}\int_0^\infty\!\! dk\;k^2\,
   \frac{M}{E(k)} \big[\, f_-(k) + f_+(k) \,\big] ,
\label{eq:nsmedT}\\[2pt]
\varepsilon^{\rm med}(\nu, M, g, T) &= \frac{g}{2\pi^2}\int_0^\infty\!\! dk\;
   k^2\, E(k) \big[\, f_-(k) + f_+(k) \,\big] ,
\label{eq:epsmedT}\\[2pt]
P^{\rm med}(\nu, M, g, T) &= \frac{g}{6\pi^2}\int_0^\infty\!\! dk\;
   \frac{k^4}{E(k)} \big[\, f_-(k) + f_+(k) \,\big] ,
\label{eq:PmedT}
\end{align}
%
and the entropy density of the species follows from the Euler relation for
one free gas,
%
\begin{equation}
s(\nu, M, g, T) \;=\;
  \frac{\varepsilon^{\rm med} + P^{\rm med} - \nu\, n^{\rm med}}{T} ,
\label{eq:sT}
\end{equation}
%
which is how the code obtains it, rather than by integrating
$-[f\ln f + (1-f)\ln(1-f)]$ directly. Note the sign structure: the
\emph{number} density is the difference of the two occupations, because an
antiparticle carries the opposite charge, while $n^s$, $\varepsilon$ and $P$
are sums. Setting $T \to 0$ at fixed $\nu > M$ returns
Eqs.~\eqref{eq:nmed}--\eqref{eq:s} exactly, with $f_- \to \theta(\kF - k)$ and
$f_+ \to 0$.

The cut-off subtraction of Sec.~\ref{sec:cutoff} is unchanged at $T > 0$:
$n^{s,{\rm vac}}$ and $\varepsilon^{\rm vac}$ depend on $(M, g, \Lambda)$
alone. The Dirac sea is a property of the vacuum, and $\Lambda$ regularizes it
there and not in the medium, so there is no temperature in either term to
carry.

\paragraph{How the integrals are evaluated, and one thing to know about it.}
Equations~\eqref{eq:nmedT}--\eqref{eq:PmedT} are not written a second time
inside this model: they come from
\texttt{eos.general.fermi\_integrals.solve\_fermi\_jel}, the single home for
the Fermi integrals of this repository (CLAUDE.md \S7), which uses the
rational approximation of Johns, Ellis and
Lattimer~\cite{JohnsEllisLattimer1996}, accurate to about $10^{-4}$
relative. That routine evaluates the \emph{exact} closed forms of
Sec.~\ref{sec:kinetic} at $T = 0$ and the fit at every $T \ne 0$, and the fit
does not converge back to the closed form: the answer steps discontinuously
the moment $T \ne 0$ and then \emph{stays} at that offset as $T$ falls. Flat
from $T = 10^{-3}$~MeV downward, the step is
%
\begin{equation}
\frac{\Delta n}{n} = \begin{cases}
  +6.9\times10^{-6} & u\ \text{at}\ \nu = 400,\ M = 5.5\ \mathrm{MeV},\\
  -6.3\times10^{-5} & s\ \text{at}\ \nu = 500,\ M = 140.7\ \mathrm{MeV},\\
  -3.0\times10^{-6} & n\ \text{at}\ \nu = 1000\ \mathrm{MeV},
\end{cases}
\label{eq:jelfloor}
\end{equation}
%
with the same numbers in $\varepsilon$. Two consequences, both deliberate.
First, ``take $T$ small and compare against the $T = 0$ answer'' is not a
validation route below about $10^{-4}$ relative, and the continuity check in
\texttt{verify} states that floor rather than chasing it. Second, $T = 0$ is
kept as an exact special case rather than routed through the fit for
smoothness: doing so would move every number this model has ever published,
including the golden values of CLAUDE.md \S12, to buy a $10^{-5}$ cosmetic
gain. The entropy itself carries no such offset --- it is zero in both
branches at $T = 0$ --- so Eq.~\eqref{eq:sT} vanishes smoothly.

\subsection{The cut-off, and the one way it enters}
\label{sec:cutoff}

The NJL interaction is not renormalizable and is regularized by a
three-momentum cut-off $\Lambda$ applied to the \emph{vacuum} (Dirac sea)
integrals of the quark sector. Baryons and leptons carry no cut-off. Writing
$y = \Lambda/M$, the quark scalar density and energy density are the medium
terms above \emph{minus} a $\kF$-independent vacuum term,
%
\begin{align}
n^s(\kF, M, g, \Lambda) &= n^{s,{\rm med}}(\kF, M, g)
  \;-\; \underbrace{\frac{g M^3}{4\pi^2}
        \Big[\, y\sqrt{y^2+1} - \operatorname{arcsinh} y \,\Big]}_{
        \textstyle n^{s,{\rm vac}}(M, g, \Lambda)} ,
\label{eq:nsvac}\\[2pt]
\varepsilon(\kF, M, g, \Lambda) &= \varepsilon^{\rm med}(\kF, M, g)
  \;-\; \underbrace{\frac{g M^4}{16\pi^2}
        \Big[\, y(2y^2+1)\sqrt{y^2+1} - \operatorname{arcsinh} y \,\Big]}_{
        \textstyle \varepsilon^{\rm vac}(M, g, \Lambda)} ,
\label{eq:epsvac}
\end{align}
%
and the number density Eq.~\eqref{eq:nmed} is untouched, the Dirac sea
carrying no net baryon number. Two properties of this split matter and are
easy to lose:

\begin{itemize}
\item \textbf{The medium integrals are not cut off.} Only the vacuum
  subtraction carries $\Lambda$. This is not a detail: the quark kinetic
  potential $\nu_q$ exceeds $\Lambda = 602.3$~MeV above
  $\nB \approx 3\;\mathrm{fm^{-3}}$, so a cut applied to the medium integral
  would truncate the physical Fermi sea and change the equation of state
  exactly where it is being used.
\item \textbf{The vacuum terms depend only on $(M, g, \Lambda)$.} They are
  additive constants at fixed constituent mass, independent of $\kF$ and ---
  this is what makes the finite-temperature extension a one-function change ---
  independent of $T$. In the code, Eqs.~\eqref{eq:nsvac} and
  \eqref{eq:epsvac} are the model's own analytic terms, added outside the
  call into \texttt{eos.general}.
\end{itemize}

\noindent
Because $n^{s,{\rm vac}} > 0$, the quark scalar density is \emph{negative} in
vacuum: at the shipped parameters and the vacuum masses of
Eq.~\eqref{eq:eps0value}, $n^s_u(\kF = 0) = n^s_d(\kF = 0) =
-1.8433\;\mathrm{fm^{-3}}$ and $n^s_s(\kF = 0) = -2.2270\;\mathrm{fm^{-3}}$.
That negative value
\emph{is} the chiral condensate, up to a positive factor, and it is what makes
the constituent masses large in vacuum through Eq.~\eqref{eq:gap}.

\section{Parameters}
\label{sec:parameters}

Every number below is an argument (\texttt{eos.enjl.Parameters}, a frozen
dataclass), never a module-level constant. Two groups: the NJL set proper,
which is the RKH parametrization of Rehberg, Klevansky and
H\"ufner~\cite{RehbergKlevanskyHufner1996}, and the density-dependent
structural functions fitted by Xia~\cite{Xia2024} (his Table~I).

\begin{center}
\begin{tabular}{llll}
\toprule
symbol & code & shipped value & meaning \\
\midrule
$\Lambda$   & \texttt{Lambda} & $602.3$~MeV & three-momentum cut-off \\
$m_{u0}, m_{d0}$ & \texttt{m\_u0, m\_d0} & $5.5$~MeV & current light-quark masses \\
$m_{s0}$    & \texttt{m\_s0}  & $140.7$~MeV & current strange-quark mass \\
$G_S$       & \texttt{GS}     & $1.835/\Lambda^2$~MeV$^{-2}$ & scalar coupling,
                                $G_S = g_\sigma^2/4m_\sigma^2$ \\
$K$         & \texttt{K}      & $12.36/\Lambda^5$~MeV$^{-5}$ &
                                't~Hooft determinant coupling~\cite{tHooft1976} \\
\midrule
$a_S, b_S, n_S$ & \texttt{aS, bS, nS} & $0.4413715$, $0.4076285$, $0.16$~fm$^{-3}$
                                & structural function $\alpha_S(\nB)$ \\
$a_V, b_V, n_V$ & \texttt{aV, bV, nV} & $3.566049$, $1.062771$, $0.214$~fm$^{-3}$
                                & $\omega$ coupling $\Gamma_\omega(\nB)$ \\
$a_{TV}, b_{TV}, n_{TV}$ & \texttt{aTV, bTV, nTV}
                                & $0.5014459$, $0.0117601$, $0.1$~fm$^{-3}$
                                & $\rho$ coupling $\Gamma_\rho(\nB)$ \\
\midrule
$f_\Lambda$ & \texttt{f\_Lambda} & $1.0626$ & fixed by $U_\Lambda(n_{\rm sat}) = -30$~MeV \\
$f_q$       & \texttt{f\_q}      & $0.5$ & quark vector rescaling; the study
                                           uses $0.5$, $0.7$, $1.0$ \\
$B$         & \texttt{B\_GeV\_fm3} & $1.0$~GeV\,fm$^{-3}$ & Pauli-blocking
                                           strength, Eq.~\eqref{eq:baryonmass} \\
\midrule
$m_e$, $m_\mu$ & \texttt{m\_e, m\_mu} & $0.511$, $105.66$~MeV & lepton masses \\
$m_\sigma, m_\omega, m_\rho$ & \texttt{m\_sigma, m\_omega, m\_rho}
            & $630$, $10^5$, $769$~MeV & bare meson masses, inert here \\
\bottomrule
\end{tabular}
\end{center}

The three density-dependent functions, with $\nB$ the total baryon density,
are
%
\begin{align}
\alpha_S(\nB)      &= a_S\, e^{-\nB/n_S} + b_S ,
\label{eq:alphaS}\\
\Gamma_\omega(\nB) &= 4 G_S \big[\, a_V\, e^{-\nB/n_V} + b_V \,\big] ,
\label{eq:Gomega}\\
\Gamma_\rho(\nB)   &= 9 \times 4 G_S
                      \big[\, a_{TV}\, e^{-\nB/n_{TV}} + b_{TV} \,\big] ,
\label{eq:Grho}
\end{align}
%
where $\Gamma_\omega \equiv g_\omega^2/m_\omega^2$ and
$\Gamma_\rho \equiv g_\rho^2/m_\rho^2$: only the ratio $g^2/m^2$ is determined
by the model, which is why the bare meson masses in the table are inert. They
are carried because the Thomas--Fermi treatment of
Ref.~\cite{XiaMaruyamaYasutakeTatsumi2024} needs them to fix the interaction
ranges once gradient terms are switched on. Note that its
$m_\omega = 10^5$~MeV is not a typo for $105$: that work deliberately adopts a
very large $\omega$ mass to suppress density fluctuations in its Thomas--Fermi
solve, and says so.

\paragraph{The factor $9$ in $\Gamma_\rho$.} Equation~\eqref{eq:Grho} carries a
factor $9$ that the printed Eq.~(22) of Ref.~\cite{Xia2024} does not. It is
required by the isospin source used here, $J_\rho = \sum_i f_i \tau_i n_i$
with $\tau_p = +1$, $\tau_n = -1$ (Eq.~\ref{eq:Jrho}), and it is confirmed
twice, independently:
%
\begin{itemize}
\item the published symmetry energies are reproduced with it and not without
  it --- $E_{\rm sym}(0.1) = 25.50$ and $E_{\rm sym}(n_{\rm sat}) = 31.55$~MeV
  measured, against $25.5$ and $31.5$ published, where the literal Eq.~(22)
  gives $13.3$ and $20.2$;
\item on the nucleonic rows of the author's own tables the isospin splitting
  of the printed potentials gives
  $g_\rho\rho = \tfrac12[(E_F^n - E_F^p) - (\mu_n - \mu_p)]$, and dividing by
  $J_\rho = n_p - n_n$ returns exactly $9.0000\times$ Eq.~(22) at every one of
  those densities, with no fit involved.
\end{itemize}
%
The $\omega$ channel needs no such factor, because its source
Eq.~\eqref{eq:Jomega} already carries $N_i = 3$ for the baryons. The
constant lives in the code as \texttt{eos.enjl.parameters.RHO\_FACTOR}.

\paragraph{Three routes to a parameter set.} Model parameters are always
arguments, so all three have to exist. \emph{By name:}
\texttt{Parameters.default()} is the set of Figs.~4--6, and
\texttt{Parameters.named(name)} takes any of the six published $(f_q, B)$
combinations --- \texttt{fq0.5\_B0}, \texttt{fq0.5\_B1}, \texttt{fq0.7\_B0},
\texttt{fq0.7\_B1}, \texttt{fq1.0\_B0}, \texttt{fq1.0\_B1} --- named after the
author's own tables, an unknown name raising \texttt{KeyError} that lists them.
\emph{A new set:} every field carries a default, so
\texttt{Parameters(f\_q=..., GS=...)} names only what changes; the dataclass is
frozen, so \texttt{dataclasses.replace} is how a set already in hand is
modified. \emph{From nuclear-matter parameters:} no route. ENJL spans both
baryons and quarks and has no \texttt{nmp.py}: its couplings are fixed by
vacuum data and by the author's tables, not by a saturation-property list.

\section{The mean field}
\label{sec:meanfield}

\subsection{Effective scalar densities and the gap equation}

The scalar source of the gap equation is not the quark scalar density alone.
A baryon is a cluster of three quarks, so its scalar density contributes to
the condensate of each of its valence flavours, weighted by the structural
function $\alpha_S$:
%
\begin{equation}
\nbs_q \;=\; \min\!\left(\,
   n^s_q(\kF^q, M_q, 6, \Lambda)
   \;+\; \alpha_S(\nB) \sum_{i = p, n, \Lambda} N^q_i\, n^s_i
   \,,\; 0 \,\right) ,
\label{eq:nbar}
\end{equation}
%
with $n^s_q$ from Eq.~\eqref{eq:nsvac} (cut-off included) and $n^s_i$ from
Eq.~\eqref{eq:nsmed} at $\Lambda = 0$. The constituent masses then satisfy the
three-flavour NJL gap
equation~\cite{NambuJonaLasinio1961,RehbergKlevanskyHufner1996} with the
't~Hooft determinant term,
%
\begin{equation}
\boxed{\;
M_u = m_{u0} - 4 G_S \nbs_u + 2 K\, \nbs_d\, \nbs_s \;}
\label{eq:gap}
\end{equation}
%
and cyclically for $d$ and $s$. Two statements about
Eqs.~\eqref{eq:nbar}--\eqref{eq:gap} are properties of \emph{this
implementation} and not of the printed equations, and both change numbers:

\begin{itemize}
\item \textbf{The determinant term is written over the other two flavours.}
  Reference~\cite{Xia2024} writes it as
  $2 K\, \nbs_u \nbs_d \nbs_s / \nbs_q$, which is the same number wherever
  $\nbs_q \ne 0$ but is $0/0$ at chiral restoration, exactly where a solver
  needs it. Equation~\eqref{eq:gap} is that expression with the removable
  singularity removed.
\item \textbf{$\nbs_q$ is capped at zero from above.} $\nbs_q$ is the
  condensate up to a positive factor ($g_\sigma\sigma_q = 4 G_S \nbs_q$), so
  it is negative in vacuum and rises towards zero as chiral symmetry is
  restored. A positive value would be a condensate of the wrong sign and
  would drive $M_q$ \emph{below} its current mass. The cap binds because the
  cluster term of Eq.~\eqref{eq:nbar} is positive and grows with baryon
  density: in symmetric nucleonic matter at $\nB = 10\;\mathrm{fm^{-3}}$ the
  uncapped expression returns $+0.105\;\mathrm{fm^{-3}}$ for the $u$ and $d$
  flavours, long after the light condensates have vanished, while the $s$
  flavour --- which the nucleons do not feed --- is still at
  $-2.072\;\mathrm{fm^{-3}}$ and is not capped. Once $\nbs_q = 0$,
  Eq.~\eqref{eq:gap} returns $M_q = m_{q0}$
  exactly and that flavour decouples from the determinant term of the other
  two.
\end{itemize}

\noindent
The gap equation is the model's only genuine self-consistency: the three
$M_q$ are the unknowns, and everything in the rest of this section is
algebraic once they are known.

\subsection{Baryon masses}

A baryon's mass is the sum of its valence constituent masses, interpolated by
$\alpha_S$ between the current and constituent values, plus a Pauli-blocking
term proportional to the baryon density carried by \emph{deconfined} quarks,
%
\begin{equation}
M_i \;=\; \sum_q N^q_i \big[\, m_{q0} + \alpha_S(\nB)\,(M_q - m_{q0}) \,\big]
      \;+\; B\, \nB^{Q} ,
\qquad
\nB^{Q} \equiv \frac{n_u + n_d + n_s}{3} .
\label{eq:baryonmass}
\end{equation}
%
The second term is what makes baryons unbind: as quarks are liberated,
$\nB^Q$ grows, every baryon mass is pushed up, and past some density the
baryons dissolve --- the Mott (deconfinement) transition. $B = 0$ switches
that mechanism off, and the two shipped values $B = 0$ and
$1\;\mathrm{GeV\,fm^{-3}}$ bracket the study. In vacuum
($\alpha_S(0) = 0.849$, $\nB^Q = 0$) Eq.~\eqref{eq:baryonmass} gives
$M_N = 938.89$ and $M_\Lambda = 1113.68$~MeV, against the published $938.9$
and $1113.7$.

\subsection{Vector fields}

Only the products $g_\omega\omega_0$ and $g_\rho\rho_0$ enter, and each is its
coupling times its source,
%
\begin{align}
J_\omega &= \sum_i f_i\, N_i\, n_i
          = 3\,(n_p + n_n + f_\Lambda n_\Lambda)
            + f_q\,(n_u + n_d + n_s) ,
\qquad N_i \equiv \sum_q N^q_i ,
\label{eq:Jomega}\\
J_\rho   &= \sum_i f_i\, \tau_i\, n_i
          = (n_p - n_n) + f_q\,(n_u - n_d) ,
\label{eq:Jrho}\\
g_\omega\omega &= \Gamma_\omega(\nB)\, J_\omega ,
\qquad
g_\rho\rho     = \Gamma_\rho(\nB)\, J_\rho .
\label{eq:fields}
\end{align}
%
$N_i = 3$ for every baryon and $N_q = 1$ for every quark: the $\omega$ couples
to quark number, and a baryon couples three times as strongly because it
contains three quarks. Leptons do not appear, $f_e = f_\mu = 0$.

\subsection{Rearrangement}

Because $\alpha_S$, $\Gamma_\omega$ and $\Gamma_\rho$ depend on $\nB$, the
variation of the energy density with respect to a density picks up terms from
the couplings themselves. These are the rearrangement self-energies, and they
enter the chemical potentials and hence the pressure, \emph{never} the energy
density:
%
\begin{align}
\SR_b &= \tfrac12\,\Gamma_\omega'(\nB)\, J_\omega^2
       + \tfrac12\,\Gamma_\rho'(\nB)\, J_\rho^2
       + \alpha_S'(\nB) \sum_{i = p,n,\Lambda}
         \Big[\sum_q N^q_i\,(M_q - m_{q0})\Big] n^s_i ,
\label{eq:SigmaRb}\\
\SR_q &= \tfrac13\, B \sum_{i = p,n,\Lambda} n^s_i
       \;+\; \tfrac13\, \SR_b ,
\label{eq:SigmaRq}
\end{align}
%
with $\Gamma' \equiv \mathrm{d}\Gamma/\mathrm{d}\nB$ and
$\alpha_S' \equiv \mathrm{d}\alpha_S/\mathrm{d}\nB$, all three read off
Eqs.~\eqref{eq:alphaS}--\eqref{eq:Grho} analytically. The asymmetry of
Eq.~\eqref{eq:SigmaRq} is worth reading twice: the Pauli-blocking term
$B\,\nB^Q$ of Eq.~\eqref{eq:baryonmass} raises \emph{baryon} masses, but
$\nB^Q$ is a \emph{quark} density, so differentiating with respect to a quark
density acts back on the quark potential --- and the factor $\tfrac13$ is
$\partial \nB^Q/\partial n_q$. The second term of Eq.~\eqref{eq:SigmaRq} is
$\tfrac13\SR_b$ for the same reason: a quark carries one third of a baryon's
worth of $\nB$, which is what the density-dependent couplings respond to.

That $\SR$ is placed correctly is not asserted --- it is measured. The
statement that fixes it is the thermodynamic definition of the chemical
potential,
%
\begin{equation}
\mu_i \;=\; \frac{\partial \varepsilon}{\partial n_i}
   \bigg|_{n_{j \ne i}} ,
\label{eq:HVH}
\end{equation}
%
which holds to $1.8\times10^{-6}$~MeV --- the floor of the central difference
--- over baryonic, hyperonic, mixed and quark-dominated compositions at
$\nB = 0.16$--$3\;\mathrm{fm^{-3}}$ and three parameter sets, on potentials of
order $10^3$~MeV. Drop either rearrangement term and Eq.~\eqref{eq:HVH} fails
by tens of MeV.

\paragraph{The one exception, and it is the cap.} Equation~\eqref{eq:HVH}
holds because the gap equation makes $\varepsilon$ stationary with respect to
each $M_q$. Where the cap of Eq.~\eqref{eq:nbar} binds, $\nbs_q$ stops
responding to the densities and that stationarity is lost in the capped
channel. It costs nothing in either of the two symmetric situations --- with
no light flavour capped nothing is clamped, and with \emph{both} $u$ and $d$
capped the determinant term vanishes in both light channels, $M_u = M_d =
m_{q0}$ exactly, and the state sits in a flat region --- so the only states
affected are those with exactly one light flavour at the cap. One such state
is reached in the checks: $f_q = 0.5$, $B = 0$ at
$\nB = 0.8\;\mathrm{fm^{-3}}$, where Eq.~\eqref{eq:HVH} misses by
$6.9\times10^{-2}$~MeV on $\mu_\Lambda = 1419$~MeV, i.e. $4.8\times10^{-5}$
relative. That is below the $0.05$--$0.20$~MeV at which the engine is
validated against the author's own tables
(Sec.~\ref{sec:validation}), so it does not show up there; it is recorded
because it is a property of the cap rather than of the arithmetic, and
removing it means changing how Eq.~\eqref{eq:nbar} is regularized, which is a
physics decision and not a refactor.

\subsection{Chemical potentials}

With $\nu_i = \sqrt{k_{F,i}^2 + M_i^2}$ the kinetic (effective) potential of
Sec.~\ref{sec:kinetic},
%
\begin{align}
\mu_i &= \nu_i + f_i\,\big(3\, g_\omega\omega + \tau_i\, g_\rho\rho\big)
         + \SR_b , &&i = p, n, \Lambda ,
\label{eq:mub}\\
\mu_q &= \nu_q + f_q\,\big(g_\omega\omega + \tau_q\, g_\rho\rho\big)
         + \SR_q , &&q = u, d, s ,
\label{eq:muq}\\
\mu_\ell &= \nu_\ell = \sqrt{k_{F,\ell}^2 + m_\ell^2} ,
         &&\ell = e, \mu .
\label{eq:mul}
\end{align}
%
The factor $3$ on $g_\omega\omega$ in Eq.~\eqref{eq:mub} and its absence in
Eq.~\eqref{eq:muq} is the same three-valence-quark counting as in
Eq.~\eqref{eq:Jomega}; the $\rho$ term carries no such factor because
$\tau_i$ already distinguishes the two nucleons.

\section{Energy density, pressure, entropy}
\label{sec:assembly}

The energy density sums the single-species kinetic terms of
Sec.~\ref{sec:kinetic} --- with the cut-off subtraction on the quarks and
without it on baryons and leptons --- and adds the condensate and vector
terms of the interaction:
%
\begin{equation}
\boxed{\;
\varepsilon \;=\; \sum_i \varepsilon(\nu_i, M_i, g_i, \Lambda_i, T)
  \;+\; 2 G_S \sum_q \big(\nbs_q\big)^2
  \;-\; 4 K\, \nbs_u \nbs_d \nbs_s
  \;+\; \tfrac12\,\Gamma_\omega J_\omega^2
  \;+\; \tfrac12\,\Gamma_\rho   J_\rho^2
  \;+\; \varepsilon_\gamma + \varepsilon_\nu
  \;-\; \varepsilon_0 \;}
\label{eq:eps}
\end{equation}
%
with $\Lambda_i = \Lambda$ for $i \in \{u,d,s\}$ and $\Lambda_i = 0$
otherwise, and $\varepsilon_\gamma$, $\varepsilon_\nu$ the optional thermal
sectors of Sec.~\ref{sec:thermalsectors}, zero unless asked for and
identically zero at $T = 0$. \textbf{No rearrangement term appears in
Eq.~\eqref{eq:eps}}, and that is the invariant a density-dependent mean field
lives or dies by: $\SR$ is the derivative of the couplings, so it belongs to
$\mu$ and to $P$ and putting it in $\varepsilon$ would count it twice
(CLAUDE.md \S8; checked by \texttt{eos/enjl/verify}).

The entropy density is the sum of the single-species entropies of
Eq.~\eqref{eq:sT}, and of the thermal sectors where they are on:
%
\begin{equation}
\boxed{\; s \;=\; \sum_i s(\nu_i, M_i, g_i, T) \;+\; s_\gamma + s_\nu \;}
\label{eq:stot}
\end{equation}
%
\textbf{and nothing else in the model carries entropy.} The Dirac-sea
subtraction is a vacuum term, and the condensate and vector terms of
Eq.~\eqref{eq:eps} are functions of the densities alone, so neither has a
temperature derivative at fixed composition. At $T = 0$ every term of
Eq.~\eqref{eq:stot} is exactly zero, by Eq.~\eqref{eq:s} --- not to a
tolerance, but identically, which is what keeps the frozen $T = 0$ numbers of
CLAUDE.md \S12 bit-for-bit across the finite-temperature extension.

The pressure is the Euler relation,
%
\begin{equation}
P \;=\; T s + \sum_i \mu_i n_i - \varepsilon ,
\label{eq:P}
\end{equation}
%
and it is how the code computes $P$ --- not by summing Eq.~\eqref{eq:PmedT}
over species. The two agree, but only once the interaction and rearrangement
contributions to the pressure are added by hand, and Eq.~\eqref{eq:P} adds
them automatically through the $\SR$ inside every $\mu_i$. It also hands the
$\mu = 0$ thermal sectors their own pressure without their appearing in it,
since $\varepsilon_X + P_X = T s_X$ for those. The free energy density is
$f = \varepsilon - Ts = -P + \sum_i \mu_i n_i$, and reduces to
$f = \varepsilon$ at $T = 0$.

\subsection{The thermal sectors}
\label{sec:thermalsectors}

Two sectors carry no conserved charge, so they enter $\varepsilon$, $P$ and
$s$ and \emph{no equation of the solve}. Both are off by default and are
selected by the flags CLAUDE.md \S4 names.

\texttt{photons}: a blackbody gas at $\mu = 0$, $g = 2$,
%
\begin{equation}
P_\gamma = \frac{\pi^2}{45}\,\frac{T^4}{(\hbar c)^3} , \qquad
\varepsilon_\gamma = 3 P_\gamma , \qquad
s_\gamma = \frac{4\pi^2}{45}\,\frac{T^3}{(\hbar c)^3} .
\label{eq:photon}
\end{equation}

\texttt{thermal\_neutrinos}: the neutrino flavours the mode does \emph{not}
track, each as a massless $\mu = 0$ gas with $g = 1$ and its antiparticle, so
that per flavour
%
\begin{equation}
\varepsilon_{\nu,1} = \frac{7}{8}\cdot\frac{\pi^2}{15}\,
   \frac{T^4}{(\hbar c)^3} , \qquad
P_{\nu,1} = \tfrac13 \varepsilon_{\nu,1} , \qquad
s_{\nu,1} = \frac{4}{3}\,\frac{\varepsilon_{\nu,1}}{T} ,
\label{eq:thermalnu}
\end{equation}
%
and $\varepsilon_\nu = N_\nu\,\varepsilon_{\nu,1}$. \textbf{The count $N_\nu$
is a property of the mode.} This model carries the electron neutrino in the
composition, and only where a mode holds $Y_{L_e}$; the muon family is
transparent here ($\mu_\mu = \mu_e - \mu_{\nu_e}$) and the tau family is never
carried. So $N_\nu = 3$ in \texttt{beta\_eq\_neutrinoless} and in the
fixed-fraction modes, and $N_\nu = 2$ in
\texttt{beta\_eq\_neutrino\_trapped}, where $\nu_e$ is already an unknown.

\paragraph{The vacuum constant $\varepsilon_0$.} Equation~\eqref{eq:eps}
without $\varepsilon_0$ does not vanish in vacuum: the cut-off quark sea
carries a large negative energy. $\varepsilon_0$ is Eq.~\eqref{eq:eps}
evaluated at zero density, so that $\varepsilon(\text{vacuum}) = 0$,
%
\begin{equation}
\varepsilon_0 \;=\; \sum_q \varepsilon\big(0, M_q^{\rm vac}, 6, \Lambda\big)
  \;+\; 2 G_S \sum_q \big(n^{s,{\rm vac}}_q\big)^2
  \;-\; 4 K \prod_q n^{s,{\rm vac}}_q ,
\label{eq:eps0}
\end{equation}
%
with $M_q^{\rm vac}$ the solution of Eqs.~\eqref{eq:nbar}--\eqref{eq:gap} at
zero density (no baryons, so no cluster term). It is a property of the vacuum
and depends only on $(\Lambda, m_{q0}, G_S, K)$ --- not on $f_q$, not on $B$,
not on density, and not on temperature. Measured:
$M_u^{\rm vac} = M_d^{\rm vac} = 367.6483$~MeV,
$M_s^{\rm vac} = 549.4792$~MeV, and
%
\begin{equation}
\varepsilon_0 = -4263.8455\;\mathrm{MeV\,fm^{-3}} .
\label{eq:eps0value}
\end{equation}
%
That all five of the author's tables, spanning three values of $f_q$ and two
of $B$, return the same constant to $8\times10^{-7}$ relative is itself a
check on Eq.~\eqref{eq:eps0}.

\section{Conserved charges}
\label{sec:charges}

The conserved-charge densities are the species sums of the quantum numbers of
Sec.~\ref{sec:dof}, with leptons excluded from all three
(\texttt{eos.general.basis}):
%
\begin{align}
\nB &= n_p + n_n + n_\Lambda + \tfrac13\,(n_u + n_d + n_s) ,
\label{eq:nB}\\
\nC &= n_p + \tfrac23 n_u - \tfrac13\,(n_d + n_s) ,
\label{eq:nCdef}\\
\nS &= n_\Lambda + n_s ,
\label{eq:nSdef}
\end{align}
%
and $Y_C = \nC/\nB$, $Y_S = \nS/\nB$. Equation~\eqref{eq:nSdef} carries the
repository's sign, $S = +1$ per $s$ quark, so $Y_S \ge 0$ in strange matter
here where the PDG convention would give $Y_S \le 0$.

Reference~\cite{Xia2024} does not introduce $\muC$ and $\muS$: it writes the
equilibrium condition directly in the physical charge $q_i$ and the electron
potential, Eq.~\eqref{eq:beta} below. The translation is exact and worth
stating, because it is what an adapter into a composite engine would use.
Applying $\mu_i = B_i\muB + C_i\muC + S_i\muS$ to $p$ and $n$ gives
$\muC = \mu_p - \mu_n$, and comparing with Eq.~\eqref{eq:beta} gives
%
\begin{equation}
\muB = \mu_b , \qquad \muC = -\mu_e , \qquad \muS = 0 .
\label{eq:basis}
\end{equation}
%
So this repository's beta-equilibrium condition $\muC + \mu_e = 0$ is
Eq.~\eqref{eq:beta} read in the conserved-charge basis, and $\muS = 0$ is the
statement that strangeness is not conserved on the timescale of weak
equilibrium --- $\mu_\Lambda = \mu_n$ and $\mu_s = \mu_d$ are both consequences
of it, not extra assumptions. Equivalently, in the deconfined regime the
inverse map $\muB = \mu_u + 2\mu_d$ holds; that identity is what the author's
tables print in their \texttt{munr} column, and it is the quantity a phase
construction must equate across a boundary (Sec.~\ref{sec:numerics}).

\section{The closure, and the four modes}
\label{sec:modes}

\subsection{Fixed composition: not a mode}
\label{sec:fixedcomp}

The lowest-level entry point, \texttt{thermodynamics.thermo\_from\_n}, is
given \emph{all eight} species densities and returns the state. It solves
Eqs.~\eqref{eq:nbar}--\eqref{eq:gap} for $(M_u, M_d, M_s)$ --- the model's
internal self-consistency --- and evaluates
Secs.~\ref{sec:meanfield}--\ref{sec:assembly}. Nothing about the composition
is determined by it, so it is not one of the repository's modes and is not
named like one: it is the ``block at given densities'' of the repository's
shared vocabulary, the counterpart of \texttt{thermo\_from\_mu} in the models
that solve at fixed potentials. Figures~1--3 of Ref.~\cite{Xia2024} ---
constituent masses in symmetric matter, $E/A$ of symmetric and pure neutron
matter, the $\Lambda$ potential depth --- are all evaluated this way, and they
carry no branch ambiguity because the composition is imposed.

\subsection{The four modes, as one declaration}
\label{sec:beta}

All four modes of this repository are closed here, at any $T \ge 0$, with
$s/\nB$ accepted in place of $T$ throughout. And they are not four solvers: a
mode is
a \emph{declaration}, one binary choice per conserved charge, and the residual
assembly reads it. For each charge, either its fraction is imposed and its
conjugate potential is an unknown, or its potential is set by an equilibrium
relation and the fraction comes out.

\begin{center}
\begin{tabular}{llll}
\toprule
mode & independent variables & extra unknowns & the charge row \\
\midrule
\texttt{beta\_eq\_neutrinoless} & $(\nB,\ T)$ & --- &
  $\sum_i q_i n_i = 0$ \\
\texttt{beta\_eq\_neutrino\_trapped} & $(\nB,\ Y_{L_e},\ T)$
  & $\mu_{\nu_e}$ & $\sum_i q_i n_i = 0$ \\
\texttt{fixed\_YC} & $(\nB,\ Y_C,\ T)$ & --- & $\nC = Y_C \nB$ \\
\texttt{fixed\_YC\_YS} & $(\nB,\ Y_C,\ Y_S,\ T)$ & $\muS$
  & $\nC = Y_C \nB$ \\
\bottomrule
\end{tabular}
\end{center}

\noindent
The trapped mode adds the row $(n_e + n_{\nu_e})/\nB = Y_{L_e}$, and
\texttt{fixed\_YC\_YS} the row $\nS = Y_S \nB$ --- except at $T = 0$,
$Y_S = 0$, where that row is \emph{empty} and $\muS = 0$ is imposed in its
place. Both strangeness carriers of this model have $S = +1$ (the $\Lambda$
and the $s$ quark) and nothing has $S < 0$, and at $T = 0$ there are no
antiparticles, so every term of $\nS = n_\Lambda + n_s$ is non-negative and
$\nS = 0$ forces $n_\Lambda = n_s = 0$ term by term. The row
$\nS - Y_S \nB = 0$ then holds identically, for every $\muS$, and determines
nothing: the potential is a null direction of the Jacobian rather than an
unknown of the system. Imposing $\muS = 0$ chooses a representative for a
quantity the equilibrium conditions leave free, and it changes no density, no
$\varepsilon$, no $P$ and no $s$ --- only the reported $\muS$ and, through
$\mu_i = B_i \muB + C_i \muC + S_i \muS$, the potentials of the two species
that are absent. At $T > 0$ the Fermi tails populate both species and their
$S = -1$ antiparticles, the row acquires a gradient in $\muS$, and the
potential is determined in the ordinary way. In the two fixed-fraction
modes the \texttt{leptons} flag decides whether the neutralizing electrons and
muons are added; either way they enter \emph{no equation}, because $\nC$ is
already pinned by the charge row, so they are computed afterwards from
$n_e + n_\mu = \nC$ and contribute to $\varepsilon$, $P$ and $s$ alone. With
\texttt{leptons=False} the result is electrically charged strongly interacting
matter, which is what a mixed-phase construction needs for each pure phase
before imposing global neutrality.

Species potentials are the conserved-charge projection of
Sec.~\ref{sec:charges},
%
\begin{equation}
\mu_i \;=\; B_i\,\muB \;+\; C_i\,\muC \;+\; S_i\,\muS ,
\label{eq:projection}
\end{equation}
%
over the six strongly interacting species, with the lepton potentials given by
the weak-equilibrium relations $\muC + \mu_e = \mu_{\nu_e}$ and
$\mu_\mu = \mu_e - \mu_{\nu_e}$ wherever $C$ is equilibrated. In the
neutrinoless mode $\mu_{\nu_e} = 0$ and $\muS = 0$, and
Eq.~\eqref{eq:projection} reduces to the paper's Eq.~(23), which is the form
the rest of this section is written in.

\medskip
\noindent
Weak equilibrium with free-streaming neutrinos ($\mu_{\nu} = 0$) makes every
species potential a projection of two numbers, the baryon potential $\mu_b$
and the electron potential $\mu_e$,
%
\begin{equation}
\mu_i \;=\; B_i\,\mu_b \;-\; q_i\,\mu_e ,
\label{eq:beta}
\end{equation}
%
for all eight species of Sec.~\ref{sec:dof} --- $q_i$ being the physical
electric charge, leptons included. Total electric neutrality closes the
system,
%
\begin{equation}
\sum_i q_i\, n_i \;=\; 0 ,
\label{eq:neutral}
\end{equation}
%
where the sum \emph{does} run over leptons. Equations~\eqref{eq:beta} and
\eqref{eq:neutral} are the conditions to keep apart: the first is
equilibrium, the second is neutrality, and Sec.~\ref{sec:charges}'s $\nC$ is
the non-leptonic charge that neither of them fixes directly. The muon family
is carried, not merely declared: $\mu^-$ populates from
$\mu_e = m_\mu = 105.66$~MeV upward, which at the shipped parameters is
between $\nB = 0.1$ and $0.2\;\mathrm{fm^{-3}}$, and by
$\nB = 1.2\;\mathrm{fm^{-3}}$ carries $n_\mu = 0.0204\;\mathrm{fm^{-3}}$
against $n_e = 0.0338\;\mathrm{fm^{-3}}$.

\subsection{The residual, row by row}
\label{sec:residual}

The unknown vector has ten components,
%
\begin{equation}
\mathbf{x} \;=\; \big(\,
  M_u,\; M_d,\; M_s,\;
  \mu_b,\; \muC,\;
  \nB^{Q},\;
  g_\omega\omega,\; g_\rho\rho,\;
  \SR_b,\; \SR_q \,\big)
  \;\oplus\; \big(\muS\big)_{Y_S\ \rm held}
  \;\oplus\; \big(\mu_{\nu_e}\big)_{Y_{L_e}\ \rm held} ,
\label{eq:unknowns}
\end{equation}
%
ten always, plus one potential per held fraction. $\muC$ is \emph{not} in that
second group: it is a potential of the matter whichever mode is asked, and
what the declaration changes is the row that closes it. It is longer than the three-unknown gap solve of the fixed-composition
entry point for one reason: at fixed composition the fields and the
rearrangement terms are \emph{known} from the densities, whereas here the
densities are what is being solved for. Carrying $g_\omega\omega$,
$g_\rho\rho$, $\SR_b$ and $\SR_q$ as unknowns with their defining equations as
residual rows replaces a nested fixed-point iteration by four more rows of one
outer solve, and --- the reason it is done this way --- guarantees that every
quantity entering any residual comes from one and the same state.
$\nB^Q$ is carried as an unknown for the same reason: it enters the baryon
masses through Eq.~\eqref{eq:baryonmass} \emph{before} the quark densities
that define it are known.

Given $\mathbf{x}$ and the target density $\nB^{\rm target}$, the state is
built forwards: the couplings $\alpha_S$, $\Gamma_\omega$, $\Gamma_\rho$ and
their derivatives are evaluated at $\nB^{\rm target}$; the baryon masses from
Eq.~\eqref{eq:baryonmass}; then for every species, in turn,
%
\begin{equation}
\nu_i = \mu_i - \Delta_i , \qquad
\Delta_i =
\begin{cases}
 f_i\,(3 g_\omega\omega + \tau_i\, g_\rho\rho) + \SR_b , & i = p,n,\Lambda \\
 f_q\,(  g_\omega\omega + \tau_i\, g_\rho\rho) + \SR_q , & i = u,d,s \\
 0 , & i = e,\mu
\end{cases}
\label{eq:forward}
\end{equation}
%
with $\mu_i$ from Eq.~\eqref{eq:projection}, and then $n_i$, $n^s_i$,
$J_\omega$, $J_\rho$ and the two $\SR$ expressions from
Secs.~\ref{sec:kinetic} and~\ref{sec:meanfield}. The residuals, in the order
the code assembles them, are the nine that every mode carries followed by the
charge row and one row per held fraction:

\begin{center}
\begin{tabular}{clll}
\toprule
row & residual & scale it is divided by & what it imposes \\
\midrule
1--3 & $M_q - \big[m_{q0} - 4G_S\nbs_q + 2K\nbs_{q'}\nbs_{q''}\big]$
     & $100$~MeV
     & the gap equation, Eq.~\eqref{eq:gap} \\
4 & $\sum_i B_i n_i - \nB^{\rm target}$
     & $\nB^{\rm target}$
     & baryon number, Eq.~\eqref{eq:nB} \\
5 & $\nB^{Q} - \tfrac13(n_u + n_d + n_s)$
     & $\nB^{\rm target}$
     & the quark baryon density, Eq.~\eqref{eq:baryonmass} \\
6 & $g_\omega\omega - \Gamma_\omega(\nB^{\rm target})\, J_\omega$
     & $3\,\Gamma_\omega \nB^{\rm target}$
     & the $\omega$ field, Eq.~\eqref{eq:fields} \\
7 & $g_\rho\rho - \Gamma_\rho(\nB^{\rm target})\, J_\rho$
     & $\Gamma_\rho \nB^{\rm target}$
     & the $\rho$ field, Eq.~\eqref{eq:fields} \\
8 & $\SR_b - \SR_b[\mathbf{x}]$
     & $3000$~MeV
     & baryon rearrangement, Eq.~\eqref{eq:SigmaRb} \\
9 & $\SR_q - \SR_q[\mathbf{x}]$
     & $1000$~MeV
     & quark rearrangement, Eq.~\eqref{eq:SigmaRq} \\
\midrule
10 & $\sum_i q_i n_i$ \quad\emph{or}\quad $\nC - Y_C\,\nB^{\rm target}$
     & $\nB^{\rm target}$
     & neutrality, Eq.~\eqref{eq:neutral}, \emph{or} the held $Y_C$ \\
11 & $\nS - Y_S\,\nB^{\rm target}$, \emph{or} $\muS$
     & $\nB^{\rm target}$, \emph{or} $100$~MeV
     & the held $Y_S$; the second form where that row is empty
       \hfill (\texttt{fixed\_YC\_YS} only) \\
12 & $n_e + n_{\nu_e} - Y_{L_e}\,\nB^{\rm target}$
     & $\nB^{\rm target}$
     & the held $Y_{L_e}$ \hfill (trapped mode only) \\
\bottomrule
\end{tabular}
\end{center}

\noindent
Rows 1--9 are present in every mode; rows 10--12 are what the declaration of
Sec.~\ref{sec:beta} selects, and the count of unknowns follows them exactly.
The trapped neutrinos are massless and left-handed, $g = 1$, so
$n_{\nu_e} = \mu_{\nu_e}^3/6\pi^2$.

Row~4 is what licenses evaluating the density-dependent couplings at
$\nB^{\rm target}$ rather than at the state's own baryon density: at the
solution the two are equal, and using the target makes
$\alpha_S$, $\Gamma_\omega$, $\Gamma_\rho$ constants of the residual rather
than functions of the unknowns, which removes their derivatives from the
Jacobian. The third column is the scale each row is divided by before the
convergence gate is applied: the rows carry mixed units --- MeV, MeV$^3$ ---
and a norm of the raw vector would be dominated by whichever row happens to be
largest (\texttt{eos.general.solve}). Once solved, the composition is handed
back to \texttt{thermo\_from\_n}, seeded with the constituent masses just
found, so the reported state is produced by exactly one code path whichever
entry point was called.

\paragraph{What temperature changes in the residual: no row, and one map.}
\emph{Not a single row of the table above changes at $T > 0$}, and neither
does the unknown vector Eq.~\eqref{eq:unknowns}. The reason is that the
residual is written in the \emph{potentials}: rows 1--3 equate constituent
masses, rows 4--9 equate fields and rearrangement terms to their definitions,
and the last rows equate charge densities to what the mode demands. Every one
of those statements is temperature-independent as an \emph{equation}.

What changes is the map from $\nu_i$ to the densities that fill those rows.
At $T = 0$, Eq.~\eqref{eq:forward} continues
$k_{F,i} = \sqrt{\nu_i^2 - M_i^2}$ for $\nu_i > M_i$ and $0$ otherwise, and
then $n_i$ and $n^s_i$ come from the closed forms
Eqs.~\eqref{eq:nmed}--\eqref{eq:nsmed}. At $T > 0$ there is no $k_{F,i}$: $n_i$
and $n^s_i$ come from Eqs.~\eqref{eq:nmedT}--\eqref{eq:nsmedT} evaluated at
$\nu_i$ directly, and a $\nu_i$ below $M_i$ is a thermally populated state
rather than an empty one. This is the \emph{forward} direction, potentials to
densities, so it costs no more than the integral itself.

\paragraph{The one direction that inverts, and where it sits.}
\texttt{thermo\_from\_n} (Sec.~\ref{sec:fixedcomp}) goes the other way: the
densities are given and the $\nu_i$ they correspond to are wanted. At $T = 0$
that is algebra, $k_{F,i}$ from Eq.~\eqref{eq:nmed} inverted and then
$\nu_i = \sqrt{k_{F,i}^2 + M_i^2}$. At $T > 0$ it is a genuine numerical
inversion of Eq.~\eqref{eq:nmedT}, done by the shared
\texttt{eos.general.fermi\_integrals.invert\_fermi\_density}.

Where that inversion sits is the structural part. $k_{F,i}$ depends only on
$(n_i, g_i)$, so at $T = 0$ it is computed once, \emph{outside} the gap
iteration. Its $T > 0$ counterpart $\nu_i$ depends on the \emph{mass}, and
the three constituent masses are exactly what the gap equation solves for ---
so all six strongly interacting species must be re-inverted at every gap
iteration. The two leptons have fixed masses and stay outside it. Measured,
the inversion costs about $30\,\mu$s against $0.1\,\mu$s for the closed form,
so a \texttt{thermo\_from\_n} call at $T > 0$ costs of order
$6 \times 20 \times 30\,\mu\mathrm{s} \approx 4$~ms where the $T = 0$ call
costs nothing extra. That asymmetry --- forward free, inverse not --- is why
the phase-adapter surface of Sec.~\ref{sec:frommu}, which a mixed-phase
construction consumes, is the forward one.

\paragraph{Entropy per baryon in place of temperature.}
Wherever a temperature is accepted, $s/\nB$ is accepted in its place
(CLAUDE.md \S3). It is an outer one-dimensional solve for $T$: an isentrope is
an isotherm whose temperature is found, per density, from
$s(\nB, T)/\nB = (s/\nB)^{\rm target}$, which is monotone increasing in $T$
at fixed density and so a well-posed bracketed root. Every evaluation of that
bracket is a full solve of the residual above. $s/\nB = 0$ is answered
directly at $T = 0$: it is the only entropy the exact $T = 0$ branch reaches,
and no positive bracket contains it.

\subsection{The block at given potentials}
\label{sec:frommu}
\label{sec:fromMu}

One further entry point is neither a mode nor a fixed composition:
\texttt{thermo\_from\_mu}$(\muB, \muC, \muS, T)$, which solves the model's own
self-consistency --- the gap equation, the vector fields, the rearrangement
terms \emph{and the phase's own baryon density} --- and returns the block. No
leptons, no neutrality, no held fraction: those are conditions on a
\emph{system}, and this describes matter. It is the surface this
repository's phase-adapter contract is written in, and a mixed-phase
construction consumes exactly this, once per phase.

Nine unknowns, $M_u$, $M_d$, $M_s$, $\nB$, $\nB^Q$, $g_\omega\omega$,
$g_\rho\rho$, $\SR_b$ and $\SR_q$, against the nine rows of
Sec.~\ref{sec:residual} with the baryon-number row read the other way round:
the mode solve imposes $\nB$ and finds $\muB$; this imposes $\muB$ and finds
$\nB$. Which \emph{branch} is returned is decided by the starting point, since
above a first-order transition the same potentials admit more than one
solution --- so a caller that differentiates this function numerically must
make its seed a deterministic function of the arguments, or the returned
Jacobian is not the derivative of anything.

\medskip
\noindent
\textbf{Temperature.} Every mode is closed at any $T \ge 0$, and $s/\nB$ is
accepted in place of $T$ throughout (Sec.~\ref{sec:residual}); a negative $T$
raises. What still refuses a temperature is the \emph{construction} of
Sec.~\ref{sec:construction}, not the model. \textbf{Species.} The composition
of Sec.~\ref{sec:dof} is fixed by the model: the octet beyond $\Lambda$, the
$\Delta$ quartet and a thermal meson gas are absent, and switching one on
raises rather than being silently ignored. The muon family is present and on.
Photons and thermal neutrinos are the two the caller \emph{does} choose ---
they are implemented, Eqs.~\eqref{eq:photon}--\eqref{eq:thermalnu}, off by
default and identically zero at $T = 0$. A thermal meson gas is not merely unimplemented but
inapplicable: $\sigma$, $\omega$ and $\rho$ here are auxiliary fields
eliminated in favour of $g^2/m^2$, and Eqs.~\eqref{eq:Gomega}--\eqref{eq:Grho}
leave the bare masses undetermined for exactly that reason.

\section{What a solved point returns}
\label{sec:returns}

\texttt{thermo\_from\_n} returns an \texttt{EoSPoint}, in natural units
(densities in MeV$^3$, masses and potentials in MeV, $\varepsilon$ and $P$ in
MeV$^4$) with fm-based properties for the public boundary. Every field, and
where it comes from:

\begin{center}
\begin{tabular}{lll}
\toprule
field & symbol & from \\
\midrule
\texttt{n}          & $n_i$, all eight species & the request \\
\texttt{M\_q}       & $M_u, M_d, M_s$ & the gap equation, Eq.~\eqref{eq:gap} \\
\texttt{M\_b}       & $M_p, M_n, M_\Lambda$ & Eq.~\eqref{eq:baryonmass} \\
\texttt{nu}         & $\nu_i$, all eight --- the PRIMARY kinetic quantity
                    & Eq.~\eqref{eq:forward} \\
\texttt{kF}         & $k_{F,i}$, the $T = 0$ derived view of $\nu_i$
                    & $\sqrt{\nu_i^2 - M_i^2}$ \\
\texttt{n\_s}       & $n^s_i$, all six strongly interacting species
                    & Eqs.~\eqref{eq:nsmed}, \eqref{eq:nsmedT},
                      \eqref{eq:nsvac} \\
\texttt{nbar\_s}    & $\nbs_q$ & Eq.~\eqref{eq:nbar} \\
\texttt{alpha\_S}, \texttt{Gw}, \texttt{Gr}
                    & $\alpha_S, \Gamma_\omega, \Gamma_\rho$
                    & Eqs.~\eqref{eq:alphaS}--\eqref{eq:Grho} at $\nB$ \\
\texttt{J\_omega}, \texttt{J\_rho}
                    & $J_\omega, J_\rho$ & Eqs.~\eqref{eq:Jomega}--\eqref{eq:Jrho} \\
\texttt{gomega\_omega}, \texttt{grho\_rho}
                    & $g_\omega\omega, g_\rho\rho$ & Eq.~\eqref{eq:fields} \\
\texttt{SigmaR\_b}, \texttt{SigmaR\_q}
                    & $\SR_b, \SR_q$ & Eqs.~\eqref{eq:SigmaRb}--\eqref{eq:SigmaRq} \\
\texttt{mu}         & $\mu_i$, all eight & Eqs.~\eqref{eq:mub}--\eqref{eq:mul} \\
\texttt{eps}, \texttt{P} & $\varepsilon$, $P$ & Eqs.~\eqref{eq:eps}, \eqref{eq:P} \\
\texttt{s}          & $s$, zero at $T = 0$ & Eq.~\eqref{eq:stot} \\
\texttt{T}          & the temperature it was solved at & --- \\
\texttt{n\_b}, \texttt{n\_bQ} & $\nB$, $\nB^Q$
                    & Eqs.~\eqref{eq:nB}, \eqref{eq:baryonmass} \\
\texttt{n\_C}, \texttt{n\_S}  & $\nC$, $\nS$
                    & Eqs.~\eqref{eq:nCdef}--\eqref{eq:nSdef} \\
\texttt{EperB}      & $\varepsilon/\nB - 938.9$~MeV & the paper's Fig.~2 ordinate \\
\bottomrule
\end{tabular}
\end{center}

\noindent
A table row (\texttt{table.beta\_row}) is that state flattened and made
fm-based: $\nB$, $T$, $P$, $\varepsilon$, $s$, $s/\nB$, the four potentials
$\muB$, $\muC$, $\muS$, $\mu_e$, the fractions $Y_C$ and $Y_S$, the eight
species densities, the three constituent and three baryon masses, and
%
\begin{equation}
\chi \;\equiv\; \frac{\nB^Q}{\nB} ,
\label{eq:chi}
\end{equation}
%
the fraction of the baryon density carried by deconfined quarks --- the
\texttt{fq} column of the author's tables. It plays here the part the quark
volume fraction plays in a two-engine mixed phase, and carries the same name
for that reason, but it is not a volume fraction: the two phases of this model
occupy the same volume.

\paragraph{The scalar density, and why it is not the trace identity.} $n^s_i$
is returned for all six strongly interacting species and is \emph{integrated},
Eqs.~\eqref{eq:nsmed}, \eqref{eq:nsmedT} and \eqref{eq:nsvac}, not obtained
from the trace of the energy--momentum tensor. The trace form
%
\begin{equation}
\varepsilon^{\rm med} - 3 P^{\rm med} = M\, n^{s,{\rm med}}
\label{eq:trace}
\end{equation}
%
does hold here, identically and per species, for the \emph{medium} terms: it
follows from Eqs.~\eqref{eq:epsmed}--\eqref{eq:Pmed} and \eqref{eq:nsmed} term
by term, and is verified to round-off for a baryon and for a quark with the
cut-off switched off. It fails for the quarks as they are actually used, and
for the totals, and for two separate reasons. The Dirac-sea subtraction removes
$\varepsilon^{\rm vac}$ from $\varepsilon$ and $n^{s,{\rm vac}}$ from $n^s$ but
nothing from $P$, and $\varepsilon^{\rm vac} \ne M\, n^{s,{\rm vac}}$; and the
total $\varepsilon$ and $P$ of Eqs.~\eqref{eq:eps} and \eqref{eq:P} carry
condensate, vector, rearrangement and thermal terms that Eq.~\eqref{eq:trace}
knows nothing about. So $n^s = (\varepsilon - 3P)/m^*$ is not a route to $n^s$
in this model and is not used as one; where a reader of another model's
document meets that identity, this is where it stops.

A mode solve returns a \texttt{BetaPoint}, which carries the same record plus
the density it was asked at, the mode declaration it was asked under, the
conserved-charge potentials $\muB$, $\muC$ and $\muS$ it solved for, the
lepton potential $\mu_e$, and the converged unknown vector $\mathbf{x}$ ---
which is what warm-starts the next density of a sweep, and what makes the
continuation follow one branch. \texttt{eos\_response} is not implemented for this model: the
second-derivative quantities of the CompOSE list divide into those that need
$T > 0$ (the heat capacities, the thermal index, the
isothermal/adiabatic distinction) and those that need a conditioning statement
the branch structure of Sec.~\ref{sec:numerics} has not yet settled ---
$c_s^2$ and the susceptibilities $\chi_{ab} = \partial n_a/\partial\mu_b$,
since above the model's first first-order transition more than one branch
satisfies the equilibrium conditions at the same density, and differentiating
along the branch a continuation happened to reach would return a number whose
meaning depends on the direction the table was swept in. The raise message
states the first half as ``this is a $T = 0$ model'', which the \emph{modes}
are not: they are closed at any $T \ge 0$. What is $T = 0$-only is the
construction of Sec.~\ref{sec:construction}, and the response gap is the
branch statement, not the temperature.

\section{The construction: from branches to a delivered EoS}
\label{sec:construction}

A table swept by \texttt{build\_table} is a \emph{continuation}: it follows one
branch and keeps following it past a first-order transition, into the
metastable region beyond, so it may violate $\mathrm{d} P/\mathrm{d}\nB \ge 0$. That is real
physics --- mechanical instability inside a coexistence region --- and mapping
it is what a branch sweep is for. What a structure solver must be handed is the
other object, the stable equation of state, and this section states how it is
built.

\subsection{The coexistence conditions}

Two branches $\alpha$ (low density) and $\beta$ (high density) of the same
functional coexist where they share a baryon chemical potential and a pressure,
each phase being separately charge-neutral with its own lepton population ---
the $\eta = 1$, or Maxwell, construction:
\begin{align}
  \muB^{(\alpha)} &= \muB^{(\beta)} \equiv \muB^{\text{co}},
  \label{eq:coexmu} \\
  P^{(\alpha)} + P_\ell^{(\alpha)} &= P^{(\beta)} + P_\ell^{(\beta)}
  \equiv P^{\text{co}},
  \label{eq:coexP} \\
  \nC^{(\gamma)} &= n_e^{(\gamma)} + n_\mu^{(\gamma)},
  \qquad \muC^{(\gamma)} + \mu_e^{(\gamma)} = 0,
  \qquad \gamma \in \{\alpha, \beta\}.
  \label{eq:coexneutral}
\end{align}
The lepton pressure belongs in Eq.~\eqref{eq:coexP} because at $\eta = 1$ the
leptons live inside the structures whose pressures are equated; dropping it
moves $\muB^{\text{co}}$ by tens of MeV. The two phases carry
\emph{different} $\mu_e$, and that difference is precisely what $\eta = 1$
means: neither borrows charge from the other. Eq.~\eqref{eq:coexneutral} is one
equation in the one unknown $\muC^{(\gamma)}$ per phase, solved first; then
Eq.~\eqref{eq:coexP} is one equation in the one remaining unknown
$\muB^{\text{co}}$.

The two edge densities follow from the solved endpoints,
\begin{equation}
  n_{\text{lo}} = \nB^{(\alpha)}\!\left(\muB^{\text{co}}\right), \qquad
  n_{\text{hi}} = \nB^{(\beta)}\!\left(\muB^{\text{co}}\right),
  \label{eq:coexedges}
\end{equation}
and between them no single phase is stable.

\subsection{The plateau}

For $n_{\text{lo}} \le \nB \le n_{\text{hi}}$ the system separates into the two
phases at fixed $\muB^{\text{co}}$ and $P^{\text{co}}$. With $f$ the volume
fraction occupied by the high-density phase, baryon number gives the lever rule
\begin{equation}
  \nB = (1 - f)\, n_{\text{lo}} + f\, n_{\text{hi}}
  \quad\Longrightarrow\quad
  f = \frac{\nB - n_{\text{lo}}}{n_{\text{hi}} - n_{\text{lo}}},
  \label{eq:lever}
\end{equation}
and every density averages with the same weights,
\begin{equation}
  n_i = (1 - f)\, n_i^{(\alpha)} + f\, n_i^{(\beta)},
  \qquad
  \varepsilon = (1 - f)\, \varepsilon^{(\alpha)} + f\, \varepsilon^{(\beta)},
  \qquad
  P = P^{\text{co}},
  \label{eq:plateau}
\end{equation}
with the entropy density levering like any other density. The conserved
fractions are formed from the
\emph{averaged} densities, $Y_C = \nC/\nB$ and $Y_S = \nS/\nB$, not by
averaging the two phases' own fractions, which are relative to their own
$\nB$. The effective masses $M_q$ and $M_i$ are properties of one phase and
have no value in the mixture; the implementation returns them as
\texttt{nan} across the plateau rather than averaging them into a state that
exists nowhere. $\muC$ and $\mu_e$ are returned as \texttt{nan} for the same
reason --- at $\eta = 1$ the two phases genuinely differ in them.

\paragraph{The energy density needs no further solve.} On a neutral phase at
$T = 0$ the Euler relation of Sec.~\ref{sec:returns} collapses. Writing
$\sum_i \mu_i n_i$ over matter and leptons and using
Eq.~\eqref{eq:coexneutral},
\begin{equation}
  \sum_i \mu_i n_i
  = \muB \nB + \muC \nC + \muS \nS + \mu_e (n_e + n_\mu)
  = \muB \nB + \muS \nS,
  \label{eq:eulerneutral}
\end{equation}
because $\muC \nC + \mu_e (n_e + n_\mu) = -\mu_e \nC + \mu_e \nC = 0$. In
$\beta$ equilibrium without strangeness conservation $\muS = 0$, so
\begin{equation}
  \varepsilon = \muB \nB - P .
  \label{eq:epsneutral}
\end{equation}
Applied at $\muB^{\text{co}}$ and $P^{\text{co}}$ this makes $\varepsilon$
\emph{linear} in $\nB$ across the plateau automatically, which is what
Eq.~\eqref{eq:plateau} asserts --- the two statements agree identically, and
the slope of the plateau in the $(\nB, \varepsilon)$ plane \emph{is}
$\muB^{\text{co}}$. It also means $c_s^2 = \mathrm{d} P/\mathrm{d}\varepsilon = 0$ there
exactly.

\subsection{Which branch outside a window}

Outside every window the delivered point is the branch that minimises the
energy density at that $\nB$: at $T = 0$ and fixed $\nB$, in $\beta$
equilibrium with neutrality, the stable state is the one of lowest
$\varepsilon$. Both continuations of Sec.~\ref{sec:numerics} are swept, upward
from the low-density chirally broken side and downward from the top of the
grid, and the comparison is made point by point. No branch bookkeeping enters:
where only one branch exists the comparison is trivial, and where both do it is
the physical criterion.

\subsection{Where this lives}

Locating a transition needs both branches at once, which is a two-phase
problem, so it lives in the composite engine
(\texttt{eos.mixed.construction.enjl\_coexistences}, over
\texttt{eos.mixed.boundaries.locate\_maxwell}) rather than in the model: a
model does not import a composite engine. The located windows are handed to
\texttt{eos.enjl.table.build\_constructed\_table} as an argument, and the
assembly of Eqs.~\eqref{eq:lever}--\eqref{eq:epsneutral} is pure ENJL. A
parameter set with three branches admits three pairings and the driver sweeps
the declared list of them, merging what it finds; which pairings are realised
is a property of the parameters, not of the engine.

Only $\eta = 1$ is delivered. At $0 < \eta < 1$ both a local and a global
lepton population exist, with weights $\eta$ and $1 - \eta$, and the plateau of
Eq.~\eqref{eq:plateau} is replaced by a solved mixed system at every density;
$f(\eta)$ is measured to fall monotonically from $\eta = 1$ to $\eta = 0$, by
$3 \times 10^{-2}$~MeV/fm$^3$ at the $f_q = 0.7$, $B = 1$ chiral transition,
with no interior extremum.

\section{Numerics}
\label{sec:numerics}

\paragraph{The gap solve.} Three unknowns $(M_u, M_d, M_s)$, Powell's hybrid
method, accepted on the residual of Eq.~\eqref{eq:gap} at
$10^{-12}$~MeV, which on masses of $5$--$550$~MeV is $10^{-14}$ relative or
better. The solver's own success flag is \emph{not} the gate: at this
tolerance it routinely reports ``not making good progress'' on a converged
root, because it is being asked for more precision than the residual can
supply. The residual decides. The cold start is the vacuum solution
$(367.6, 367.6, 549.5)$~MeV, which is a poor guess once the light condensates
have collapsed --- a fixed-composition request in the chirally restored region
can fail from it and must be seeded from a neighbouring point.

\paragraph{The beta-equilibrium solve.} Ten unknowns
(Eq.~\ref{eq:unknowns}), bounded least squares
(\texttt{scipy.optimize.least\_squares}) on the scaled residuals of
Sec.~\ref{sec:residual}, accepted at the repository's common gate
$10^{-10}$. The bounds are not decoration: at the shipped parameters
$\mu_b = 6.24$~GeV and $g_\omega\omega = 1.68$~GeV already at
$\nB = 3.8\;\mathrm{fm^{-3}}$, and both keep growing roughly linearly, so a
box calibrated at saturation would exclude the solution entirely above a few
times $n_{\rm sat}$; it widens with density instead, while the constituent
masses have a genuine
density-independent ceiling (their vacuum values) and $\nB^Q$ a genuine one
($\nB$ itself). Starting points are tried in order of decreasing
plausibility: the previous point of a sweep first, then the same state with
the light condensates switched off --- which is where the chirally restored
branch sits, and is what carries a sweep \emph{across} a chiral transition,
since a first-order transition puts the next root several hundred MeV away in
the quark masses. Two parameter-free cold starts, one nucleonic and one
deconfined, are available only when there is no previous point. On the shipped
grid the achieved scaled residual is at most $8.8\times10^{-13}$, median
$4.0\times10^{-15}$.

\paragraph{Continuation, and the open branch question.} A table is a
continuation, not a phase diagram. Each point is warm-started from its
neighbour, so the sequence follows one branch of the model and keeps
following it past a first-order transition into the metastable region beyond.
Cold starts are allowed only until the branch is established; permitting them
mid-sweep lets the sequence hop between branches from one density to the
next, which shows up as an equation of state that oscillates rather than one
that has a transition in it. The direction of the sweep therefore selects the
branch: upward from the low-density chirally broken side, or downward from a
deconfined guess at the top of the grid. Where only one branch exists the two
agree; where several do, they differ, and \emph{that difference is the branch
structure}.

Choosing between branches is a Maxwell construction --- at fixed $\mu_b$,
take the branch of larger $P$ --- and it cannot be done from a single sweep,
because it needs both branches at once. It is not implemented here. Three
transitions can each be first order or continuous depending on $(f_q, B)$:
the \emph{quarkyonic} onset~\cite{McLerranPisarski2007}, where quasi-free
quarks appear alongside baryons; the \emph{chiral} transition, where
$\nbs_q \to 0$ and $M_q \to m_{q0}$; and \emph{deconfinement}, where
Eq.~\eqref{eq:baryonmass}'s Pauli-blocking term unbinds the baryons. Two of
the author's own tables retain a step with $\mathrm{d}P/\mathrm{d}\nB < 0$
rather than the coexistence plateau that would replace it, which is the same
statement seen from the other side: a raw branch may violate mechanical
stability, and a construction --- not the branch map --- is what resolves it
before a table reaches a structure solver (CLAUDE.md \S8). The baryon
potential to equate across such a boundary is $\muB$ of Eq.~\eqref{eq:basis},
which in the deconfined phase is $\mu_u + 2\mu_d$ and \emph{not} the vanishing
neutron's own $\mu_n$.

\paragraph{Where finite temperature entered, and where it did not.}
Nowhere except Sec.~\ref{sec:kinetic}, as the structure of the model
predicted. The couplings are functions of $\nB$; the gap equation is the same
equation with $T$-dependent scalar densities; the rearrangement terms have the
same expressions; $\varepsilon_0$ is the same number; and \emph{not one row}
of the residual table of Sec.~\ref{sec:residual} changed, because those rows
are written in the potentials. Only
Eqs.~\eqref{eq:nmed}--\eqref{eq:s} became
Eqs.~\eqref{eq:nmedT}--\eqref{eq:sT} --- and the vacuum terms
\eqref{eq:nsvac}--\eqref{eq:epsvac} did not change even there, being
$T$-independent by construction. The pressure is the full Euler form
$P = Ts + \sum_i \mu_i n_i - \varepsilon$ of Eq.~\eqref{eq:P}.

\paragraph{What the construction still refuses.} The branch \emph{map} of
Sec.~\ref{sec:construction} --- a warm-started continuation --- runs at any
temperature. Replacing a first-order window by its plateau does not. Locating
a coexistence at $T > 0$ means equating the Gibbs free energies of the two
branches rather than $P$ and $\muB$ alone, so the entropy enters the
coexistence bookkeeping, and the lever rule of Eq.~\eqref{eq:plateau} then
averages it across the window as well. Both
\texttt{eos.enjl.table.build\_constructed\_table} and
\texttt{eos.mixed.construction} raise above $T = 0$ and say so.

\section{What the implementation reproduces, measured}
\label{sec:validation}

\subsection{Against the published numbers}

Symmetric nuclear matter is solved at fixed composition and the saturation
properties read off by central differences; the $\Lambda$ potential is
$U_\Lambda = \mu_\Lambda - M_\Lambda^{\rm vac}$ at vanishing $\Lambda$
density.

\begin{center}
\begin{tabular}{llll}
\toprule
quantity & this implementation & Ref.~\cite{Xia2024} & \\
\midrule
$n_{\rm sat}$              & $0.158297$~fm$^{-3}$ & $0.158$ & Table~II \\
$E/A$ at $n_{\rm sat}$     & $-16.010$~MeV        & $-16.0$ & Table~II \\
$K_{\rm sat}$              & $234.20$~MeV         & $234.5$ & Table~II \\
$E_{\rm sym}(n_{\rm sat})$ & $31.549$~MeV         & $31.5$  & Table~II \\
$L_{\rm sym}$              & $42.35$~MeV          & $42.4$  & Table~II \\
$E_{\rm sym}(0.1)$         & $25.500$~MeV         & $25.5$  & Sec.~II \\
$E/\nB$ at $0.1$           & $924.79$~MeV         & $924.8$ & Sec.~II \\
$E/\nB$ at $0.158$         & $922.89$~MeV         & $922.9$ & Sec.~II \\
$M_N$ at $0.16$            & $519.23$~MeV         & $519.2$ & Sec.~II \\
$M_N$, $M_\Lambda$ in vacuum & $938.89$, $1113.68$~MeV & $938.9$, $1113.7$ & Sec.~II \\
$M_u = M_d$, $M_s$ in vacuum & $367.648$, $549.479$~MeV & $367.6$, $549.5$ & Sec.~II \\
$\alpha_S(0)$, $\alpha_S(n_S)$ & $0.849000$, $0.570000$ & $0.849$, $0.57000$ & Table~I \\
$U_\Lambda(n_{\rm sat})$   & $-30.020$~MeV        & $-30$ (fitted) & Sec.~II \\
\bottomrule
\end{tabular}
\end{center}

\subsection{Against the author's own tables}

Five beta-equilibrium tables produced by the author's Maple worksheet, at
$(f_q, B) \in \{(1.0, 0), (1.0, 1), (0.7, 0), (0.7, 1), (0.5, 1)\}$ and
$\nB = 0.01$--$10\;\mathrm{fm^{-3}}$, are the tightest available constraint:
they come from the code that made the paper's figures and so pin the model far
beyond the two or three significant figures the paper prints. Reading them
correctly requires three things that are not in their column headers, and each
of them costs hundreds of MeV if missed. The scalar-density column named
\texttt{nsq} is the \emph{medium} part of Eq.~\eqref{eq:nsvac}, while the
column named \texttt{Sigmaq} is the full $\nbs_q$ of Eq.~\eqref{eq:nbar} ---
so the gap equation must be fed \texttt{Sigmaq}. The baryon potential is the
column \texttt{munr}, not \texttt{mun}: once baryons dissolve the two part
company, \texttt{munr} continuing as $\mu_u + 2\mu_d$
(Eq.~\ref{eq:basis}). And Eq.~\eqref{eq:beta} holds only for species that are
\emph{present}; a species below threshold has a stale printed potential, and
$\mu_e$ must be recovered as $\mu_d - \mu_u$ where no electrons remain.
Read that way, and with the interpolated mixed-phase rows of the $f_q = 0.5$
file and nine non-converged rows of the $f_q = 0.7$, $B = 0$ file excluded,
the engine reproduces the tables' own identities --- the gap equation to
$2\times10^{-5}$ relative, Eq.~\eqref{eq:P} to $10^{-8}$--$10^{-4}$ relative
depending on the file, and the mean-field rebuild of every printed $\mu_i$ to
between $0.005$ and $0.20$~MeV on potentials of $1000$--$2500$~MeV. These
comparisons and the per-file tolerances they need live in \texttt{test/enjl}.

\subsection{Against \texttt{eos/general}}

The free-gas parts of Eqs.~\eqref{eq:nmed}--\eqref{eq:Pmed} are the shared
$T = 0$ Fermi integrals of \texttt{eos.general.fermi\_integrals}, and the model
calls them rather than carrying a second copy. Before that routing was made,
the local closed forms and the shared ones were compared over
$g \in \{2, 6\}$, $M \in \{5.5, 140.7, 200, 300, 367.6, 500, 549.5, 938.9,
939, 1115.7\}$~MeV and $\kF \in \{1, 10, 50, 200, 500, 1200, 2200, 5000\}$~MeV,
at $\Lambda = 0$. Worst relative deviation by $\kF$:

\begin{center}
\begin{tabular}{lcccc}
\toprule
$\kF$ [MeV] & $n$ & $\varepsilon$ & $n^s$ & $P$ \\
\midrule
$1$    & $3.5\times10^{-10}$ & $6.5\times10^{-8}$  & $2.6\times10^{-7}$  & $4.1\times10^{-1}$ \\
$10$   & $1.8\times10^{-12}$ & $2.5\times10^{-11}$ & $9.3\times10^{-11}$ & $1.3\times10^{-6}$ \\
$50$   & $1.4\times10^{-13}$ & $3.3\times10^{-13}$ & $1.2\times10^{-12}$ & $5.7\times10^{-10}$ \\
$200$  & $8.5\times10^{-15}$ & $1.8\times10^{-14}$ & $3.1\times10^{-14}$ & $9.7\times10^{-13}$ \\
$\ge 500$ & $1.4\times10^{-15}$ & $1.4\times10^{-15}$ & $3.3\times10^{-15}$ & $1.4\times10^{-14}$ \\
\bottomrule
\end{tabular}
\end{center}

\noindent
The two are the same analytic expression, so the table is a floating-point
statement and not a physics one. It is quoted rather than asserted because the
$\kF = 1$~MeV entry looks alarming and is not: there $x = \kF/M \sim 10^{-3}$
and the pressure is a cancellation of two nearly equal terms in both forms.
At that point the asymptotic $P \to g\kF^5/(30\pi^2 M)$ gives
$7.880\times10^{-13}\;\mathrm{MeV\,fm^{-3}}$, this model's
$\nu n - \varepsilon$ gives $7.875\times10^{-13}$ and the shared closed form
gives $4.671\times10^{-13}$; the shared form is the one losing digits, and the
species carrying it has a density of $4\times10^{-9}\;\mathrm{fm^{-3}}$.
Everywhere a species is populated enough to matter, the agreement is
$10^{-12}$ or better, and routing the model through the shared integrals moved
no frozen quantity of \texttt{test/baseline} by more than
$1.5\times10^{-12}$ relative, against a gate of $10^{-10}$.

\subsection{Internal invariants}

\texttt{eos/enjl/verify/run\_full\_check.py} asserts, over a grid of
compositions, densities, modes and parameter sets: the Euler relation and
$f = \varepsilon - Ts = -P + \sum_i \mu_i n_i$;
$\mu_i = \partial\varepsilon/\partial n_i$, Eq.~\eqref{eq:HVH}, which is the
statement that $\SR$ is in $\mu$ and $P$ and not in $\varepsilon$, gated at
$10^{-3}$~MeV where the functional is smooth and at $10^{-1}$~MeV at the cap,
whose exception is named rather than absorbed; the gap equation satisfied at
the returned state, cap included; the charge densities of
Eqs.~\eqref{eq:nB}--\eqref{eq:nSdef} against \texttt{eos.general.basis};
Eq.~\eqref{eq:beta} for every species that is present and
Eq.~\eqref{eq:neutral} to $9\times10^{-13}$~MeV; that each fixed-fraction
mode returns the fraction it was asked for, to $7\times10^{-15}$; that
\texttt{fixed\_YC\_YS} at $Y_C = 0.5$, $Y_S = 0$ with the leptons off IS
nucleonic symmetric matter wherever its own solution is quark-free, to
$9\times10^{-12}$; that the trapped mode hits $Y_{L_e}$ and satisfies
$\muC + \mu_e = \mu_{\nu_e}$; that \texttt{thermo\_from\_mu} returns, cold,
the state whose potentials it was given, to $2\times10^{-13}$; together with
the
conserved-charge reading $\muB = \mu_u + 2\mu_d$, $\muC = -\mu_e$, $\muS = 0$
of Eq.~\eqref{eq:basis}; $0 \le c_s^2 \le 1$ on the stable stretch of a
continuation, which is LOCATED by walking the sweep rather than assumed, since
a raw branch may legitimately violate $\mathrm{d}P/\mathrm{d}\nB \ge 0$; the
vacuum masses and $\varepsilon_0$, including that $\varepsilon_0$ is identical
across parameter sets; that each mode, temperature, species flag and response
function of Sec.~\ref{sec:modes} that is said to raise, does; and that a
request the solver cannot reach comes back as a status rather than an
exception.

\bibliography{../../docs/eos}

\end{document}
